\documentclass[UTF8]{ctexart}
\usepackage[margin=2cm]{geometry}
\usepackage{amsmath,amssymb,bm,mathtools}
\usepackage{tabularx}
\usepackage{booktabs}

\title{为什么 $\mu_1 > \gamma / \alpha_H$ 让交叉项 $\tilde{S}_d(1_N \otimes v)$ 指数衰减，从而满足 Lemma 1 的 $F_2(t) \to 0$}
\author{}
\date{}

\begin{document}

\maketitle

\begin{quote}
论文：He Cai, Frank L. Lewis, Guoqiang Hu, Jie Huang,
\emph{``The adaptive distributed observer approach to the cooperative output regulation of linear multi-agent systems''}, Automatica, 2016.

本文只解释\textbf{一个}问题：条件 $\mu_1 > \gamma / \alpha_H$ 如何使交叉项 $\tilde{S}_d(1_N \otimes v)$ 指数衰减到 $0$，从而正好满足 Lemma 1 收敛证明中对 $F_2(t) \to 0$（指数）的要求。关于 $\mu_2 > \gamma / \alpha_H$ 如何让齐次矩阵 Hurwitz 的讨论，见文末附记，不构成本文主线。
\end{quote}

\hrule
\vspace{4mm}

\section{问题：我们要解释什么}

在 Lemma 2(ii) 的收敛证明中，$\tilde{\eta} = \mathrm{col}(\tilde{\eta}_1,\dots,\tilde{\eta}_N)$ 满足
\begin{equation}
\dot{\tilde{\eta}} = [(I_N \otimes S) - \mu_2(H \otimes I_q)] \tilde{\eta} + \tilde{S}_d \tilde{\eta} + \tilde{S}_d (1_N \otimes v). \tag{11}
\end{equation}

论文把 (11) 套用 Lemma 1，即系统 $\dot{x} = \epsilon F x + F_1(t)x + F_2(t)$）。要由 Lemma 1 推出 $\tilde{\eta}(t) \to 0$ \textbf{指数地}，必须保证两个扰动项都\textbf{指数衰减到 0}：
\begin{itemize}
  \item $F_1(t) = \tilde{S}_d(t)$ $\to$ 由 Lemma 2(i) 直接得到（$\tilde{S}_i \to 0$ 指数），\textbf{不需要} $\mu_1 > \gamma/\alpha_H$；
  \item \textbf{$F_2(t) = \tilde{S}_d(t)(1_N \otimes v(t))$ $\to$ 这正是要解释的交叉项}，要用统一的速率估计保证它指数衰减，需取 $\mu_1 > \gamma/\alpha_H$。
\end{itemize}

本文要证明的，就是下面这句话：
\begin{quote}
若 $\mu_1 \alpha_H - \gamma > 0$，即 $\mu_1 > \gamma/\alpha_H$，则交叉项 $\tilde{S}_d(t)(1_N \otimes v(t))$ 指数衰减到 $0$，从而满足 Lemma 1 的 $F_2(t)$ 条件。
\end{quote}

这里的“不等式等价”只指 $\mu_1 \alpha_H - \gamma > 0 \Longleftrightarrow \mu_1 > \gamma/\alpha_H$。作为对所有初值的统一充分条件，它正好解释原文为什么这样选 $\mu_1$；若初值有特殊消零结构（例如 $v(0)=0$ 或某些模态未被激发），该条件未必是某个具体轨迹的必要条件。

注意 $\mu_1$ \textbf{不出现在}齐次矩阵 $(I_N \otimes S) - \mu_2(H \otimes I_q)$ 中；它只通过 $\tilde{S}_d$ 的衰减速率影响交叉项。所以 $\mu_1 > \gamma/\alpha_H$ 与 $\mu_2 > \gamma/\alpha_H$ 职责不同（见 \S 5 与附记）。

\section{相关定义与已知估计}

先把后面要用到的符号和\textbf{已证}估计列清楚。

\subsection{Leader / exosystem 与 $\gamma$}

Leader 信号（式 (1)）：
\[
\dot{v} = S v, \qquad v(0) \text{ 给定}.
\]

定义 leader 矩阵 $S$ 的\textbf{最大实部增长率}（Lemma 2(ii)）：
\[
\gamma = \max\bigl(\Re(\sigma(S))\bigr).
\]

由常系数线性系统的基本性质，$v(t) = e^{St}v(0)$ 的范数满足
\[
\|v(t)\| \le c \cdot e^{\gamma t} \cdot (1+t^{\,m-1}),
\]
其中 $m$ 可取为相关 Jordan 块的最大尺寸；多项式因子来自非半单（non-semi-simple）Jordan 块。若 $S$ 半单（例如所有 Jordan 块都是 $1\times 1$），则简化为 $\|v(t)\| \le c \cdot e^{\gamma t}$。本文去掉了 Cai \& Huang (2016) 的“半单且零实部”假设，因此要保留这个多项式因子。

于是沿全体 follower 堆叠的信号满足
\begin{equation}
\|1_N \otimes v(t)\| \le c' \cdot e^{\gamma t} \cdot (1+t^{\,m-1})
\quad \bigl(\text{半单时即 } c' \cdot e^{\gamma t}\bigr). \tag{v-bound}
\end{equation}

即 $1_N \otimes v$ 的指数增长率至多为 $\gamma$，并可能带一个多项式因子。$\gamma > 0$ 时存在正指数增长模态；$\gamma = 0$ 且 $S$ 非半单时可能出现 ramp 等多项式增长；$\gamma < 0$ 时则指数衰减。

\subsection{分布式矩阵估计误差 $\tilde{S}_i$ 与 Lemma 2(i)}

定义每个 follower 对 $S$ 的估计误差：
\[
\tilde{S}_i = S_i - S.
\]

由 (5a) $\dot{S}_i = \mu_1 \sum_{j=0}^N a_{ij}(S_j - S_i)$，把 $\tilde{S} = \mathrm{col}(\tilde{S}_1,\dots,\tilde{S}_N)$ 堆叠，可得紧致形式：
\[
\dot{\tilde{S}} = -\mu_1 (H \otimes I_q) \tilde{S},
\qquad
H = L + \operatorname{diag}(a_{10},\dots,a_{N0}).
\]

其中 $L$ 是 follower 子图 Laplacian，$H$ 因 Assumption 4（含 leader 为根的有向生成树）而\textbf{所有特征值实部为正}引 Hu \& Hong 2007 Lemma 4）。定义图连通性 / 谱隙
\begin{equation}
\alpha_H = \min\bigl(\Re(\sigma(H))\bigr) > 0. \tag{alphaH}
\end{equation}

Lemma 2(i)给出逐块估计：对某个 $\beta_H > 0$，
\begin{equation}
\|\tilde{S}_i(t)\| \le \beta_H \|\tilde{S}_i(0)\| \, e^{-\mu_1 \alpha_H t}, \qquad i=1,\dots,N. \tag{S-tilde-decay}
\end{equation}

即\textbf{矩阵估计误差以速率 $\mu_1 \alpha_H$ 指数衰减}——这是关键：$\mu_1$ 越大、图越连通（$\alpha_H$ 越大），衰减越快。

记 $\tilde{S}_d = \mathrm{blockdiag}\{\tilde{S}_1,\dots,\tilde{S}_N\}$，由 (S-tilde-decay) 还有整体界
\begin{equation}
\|\tilde{S}_d(t)\| \le C \, e^{-\mu_1 \alpha_H t}. \tag{S-d-decay}
\end{equation}
（常数 $C$ 吸收 $\beta_H$ 与最大初值。）

技术上，如果不直接采用原文 Lemma 2(i) 给出的速率估计，而是从一般非对称矩阵 $H$ 的 Jordan 形式出发，$\tilde{S}_d(t)$ 的上界也可能多出一个多项式因子，或写成任意略慢的速率 $e^{-\mu_1(\alpha_H-\varepsilon)t}$。由于后面使用的是严格不等式 $\mu_1\alpha_H>\gamma$，这些更保守的写法仍能被指数裕度吸收，不改变本文解释的下界形式。

\section{交叉项的界：把两个速率相乘}

现在对交叉项 $\tilde{S}_d(t)(1_N \otimes v(t))$ 取范数，并用次乘性 $\|AB\| \le \|A\|\cdot\|B\|$：
\begin{align}
\|\tilde{S}_d(t)(1_N \otimes v(t))\|
&\le \|\tilde{S}_d(t)\| \cdot \|1_N \otimes v(t)\| \\[2mm]
&\le \bigl(C \, e^{-\mu_1 \alpha_H t}\bigr)
   \cdot \bigl(c' \, e^{\gamma t} \, (1+t^{\,m-1})\bigr)
   && \text{代入 (S-d-decay) 与 (v-bound)} \\[2mm]
&= C c' \cdot (1+t^{\,m-1}) \cdot e^{-(\mu_1 \alpha_H - \gamma)\, t}. \tag{cross-bound}
\end{align}

这正是要的核心界：\textbf{交叉项是“一个指数衰减项”乘以“一个多项式项”}——
\begin{itemize}
  \item 指数部分的净指数为 $-(\mu_1 \alpha_H - \gamma) t$；
  \item 多项式部分 $1+t^{m-1}$ 来自 $S$ 的非半单纯性，且\textbf{完全不依赖} $\mu_1, \alpha_H, \gamma$。
\end{itemize}

\section{何时指数衰减：$\mu_1 \alpha_H - \gamma > 0 \Longleftrightarrow \mu_1 > \gamma / \alpha_H$}

判断 (cross-bound) 这个统一上界是否指数趋于 $0$，只看指数项的净指数符号：
\[
\mu_1 \alpha_H - \gamma \;\begin{cases}
> 0 & \Rightarrow e^{-(\mu_1\alpha_H-\gamma)t} \text{ 指数衰减，且压过多项式因子}; \\[1mm]
= 0 & \Rightarrow \text{统一上界只剩多项式量，一般不能保证衰减}; \\[1mm]
< 0 & \Rightarrow \text{统一上界含正指数增长，一般不能保证衰减}.
\end{cases}
\]

由于 $\alpha_H > 0$（式 (alphaH)），不等式方向在除以 $\alpha_H$ 时不变，故
\begin{equation}
\boxed{\;\mu_1 \alpha_H - \gamma > 0 \quad \Longleftrightarrow \quad \mu_1 > \frac{\gamma}{\alpha_H}\;}. \tag{core}
\end{equation}

这就是本文要说明的\textbf{核心不等式} $\mu_1 \alpha_H - \gamma > 0$。它正是 $\mu_1 > \gamma / \alpha_H$。

\subsection{多项式因子被指数项压制}

需确认 (core) 成立时，(cross-bound) 确实指数衰减，即便有 $1+t^{m-1}$。对任意指数衰减率 $\lambda \triangleq \mu_1 \alpha_H - \gamma > 0$，以及任意多项式次数 $m-1$，都有熟知的极限
\[
\lim_{t\to\infty} (1+t^{\,m-1}) e^{-\lambda t} = 0,
\]
且收敛是指数阶的：存在 $C'', \lambda' \in (0, \lambda)$ 使 $(1+t^{m-1}) e^{-\lambda t} \leq C'' e^{-\lambda' t}$。因此 (cross-bound) $\Rightarrow$
\[
\|\tilde{S}_d(t)(1_N \otimes v(t))\|
\le C''' e^{-\lambda' t} \xrightarrow[t\to\infty]{} 0
\quad\text{指数地}.
\]

换言之，\textbf{$S$ 是否半单只改变常数项与多项式尾巴，不改变“净指数为正即可保证交叉项指数衰减”这一结论}。

\subsection{两种情形}

\begin{itemize}
  \item \textbf{$\gamma \leq 0$（S 无正实部特征值）}：此时对任意 $\mu_1 > 0$，
  $\mu_1 \alpha_H - \gamma \geq \mu_1 \alpha_H > 0$，(core) \textbf{自动成立}。即便 $\gamma=0$ 时 $v(t)$ 可能有 ramp 这类多项式增长，$e^{-\mu_1\alpha_H t}$ 仍会压过该多项式因子。不需要任何显式下界——这正是第一句中 ``for any positive $\mu$'' 的由来。
  \item \textbf{$\gamma > 0$（S 含正指数增长模态）}：必须\textbf{显式选取} $\mu_1 > \gamma / \alpha_H$。否则净指数 $\mu_1 \alpha_H - \gamma \leq 0$，上述统一估计无法保证交叉项衰减，$\tilde{\eta} \to 0$ 的证明链条在此处断裂。
\end{itemize}

\section{对接 Lemma 1：$F_2(t)$ 就是该交叉项}

现在把 (11) 与 Lemma 1 的 $\dot{x} = \epsilon F x + F_1(t)x + F_2(t)$ 逐一对齐，式 (6)）：

\begin{table}[h]
\centering
\begin{tabularx}{\textwidth}{lllX}
\toprule
Lemma 1 中的量 & 在 (11) 中的对应 & 说明 \\
\midrule
$x$ & $\tilde{\eta}$ & 待证趋于 $0$ 的误差 \\
$\epsilon F$ & $(I_N \otimes S) - \mu_2(H \otimes I_q)$ & 齐次部分，需 Hurwitz（由 $\mu_2 > \gamma/\alpha_H$ 保证，见附记） \\
$F_1(t)$ & $\tilde{S}_d(t)$ & 需 $\to 0$ 指数 \\
\textbf{$F_2(t)$} & \textbf{$\tilde{S}_d(t)(1_N \otimes v(t))$} & 需 $\to 0$ 指数 \\
\bottomrule
\end{tabularx}
\end{table}

\begin{itemize}
  \item \textbf{$F_1(t) = \tilde{S}_d \to 0$ 指数}：直接来自 Lemma 2(i) 的 (S-tilde-decay)，\textbf{与 $\mu_1 > \gamma/\alpha_H$ 无关}（对任意 $\mu_1 > 0$ 都成立。所以 $F_1$ 这一项不带来 $\mu_1$ 的下界。
  \item \textbf{$F_2(t) = \tilde{S}_d(1_N \otimes v) \to 0$ 指数}：正是 \S 3--\S 4 证明的结论；要用该统一速率估计保证它指数衰减，只需 $\mu_1 \alpha_H - \gamma > 0$，即 $\mu_1 > \gamma / \alpha_H$。
\end{itemize}

Lemma 1(i) 的结论是：若 $F$ Hurwitz 且 $F_1(t), F_2(t) \to 0$（指数），则对任意初值 $x(t) \to 0$ 指数地。于是只要
\begin{enumerate}
  \item $\mu_2 > \gamma/\alpha_H$（使 $\epsilon F$ Hurwitz），且
  \item $\mu_1 > \gamma/\alpha_H$（使 $F_2(t)$ 指数衰减）；
\end{enumerate}
两项都满足，Lemma 1 即推出 $\tilde{\eta}(t) \to 0$ 指数地，证明收口。

\begin{quote}
关键区分：$\mu_2$ 决定齐次矩阵是否 Hurwitz；$\mu_1$ 决定交叉项 $F_2$ 是否指数衰减。二者下界形式相同（$> \gamma/\alpha_H$），但职责不同。本文主线只解释 $\mu_1$ 那一侧。
\end{quote}

\section{直觉总结}

可以把交叉项理解为\textbf{两个相反趋势的乘积}：
\begin{itemize}
  \item leader 信号 $1_N \otimes v$ 的指数增长率由 $\gamma$ 控制，并可能带多项式因子（ramp 对应 $\gamma=0$ 但有多项式增长）；
  \item 矩阵估计误差 $\tilde{S}_d$ 按速率 $\mu_1 \alpha_H$ \textbf{衰减}，其中 $\mu_1$ 是耦合增益、$\alpha_H$ 是图连通性（谱隙）。
\end{itemize}

要用统一估计保证乘积 $\tilde{S}_d(1_N \otimes v)$ 归零，必须让\textbf{指数衰减快过指数增长}：
\[
\underbrace{\mu_1 \alpha_H}_{\text{衰减速率}} \;>\; \underbrace{\gamma}_{\text{增长速率}}
\quad\Longleftrightarrow\quad
\mu_1 > \frac{\gamma}{\alpha_H}.
\]

\begin{itemize}
  \item 图越连通（$\alpha_H$ 越大）$\to$ 衰减越快，所需 $\mu_1$ 越小；
  \item leader 信号增长越快（$\gamma$ 越大）$\to$ 所需 $\mu_1$ 越大；
  \item 若 leader 无正指数增长（$\gamma \leq 0$），则任意正 $\mu_1$ 都够，因为衰减侧恒为正，并会压过可能出现的多项式增长。
\end{itemize}

一句话：\textbf{$\mu_1 > \gamma/\alpha_H$ 的物理解读是——分布式观测器估计 $S$ 的收敛速度，必须胜过 leader 信号 $v$ 的增长速度，这样残差 $\tilde{S}_d$ 与 $v$ 的乘积才会指数消失，Lemma 1 的 $F_2(t) \to 0$ 条件才满足。}

\section{附记：关于 $\mu_2 > \gamma / \alpha_H$（对照，非本文主线）}

$\mu_2$ 的下界来自齐次矩阵
\[
M = (I_N \otimes S) - \mu_2 (H \otimes I_q)
\]
需为 Hurwitz。
由于 $(I_N \otimes S)$ 与 $(H \otimes I_q)$ 作用于不同张量腿而可交换，可同时三角化，$M$ 的特征值为
\[
\sigma(M) = \{\, s_k - \mu_2 h_j : s_k \in \sigma(S),\ h_j \in \sigma(H) \,\},
\]
其实部最大者为 $\gamma - \mu_2 \alpha_H$。$M$ Hurwitz $\Leftrightarrow$ $\gamma - \mu_2 \alpha_H < 0$ $\Leftrightarrow$ $\mu_2 > \gamma/\alpha_H$（由 $\alpha_H > 0$）。这与 $\mu_1$ 的下界\textbf{形式相同但职责不同}：$\mu_2$ 管齐次部分稳定，$\mu_1$ 管交叉项 $F_2$ 衰减。

Remark 4 还给出对应的 $\mu_3$ 下界：
\[
\mu_3 \ge \frac{\mu_1 \alpha_H}{\min\sigma(Q_i^\top Q_i)}
> \frac{\gamma}{\min\sigma(Q_i^\top Q_i)},
\]
其中 $Q_i$ 为 $\eta_i$ 动力学 (3) 的可观测性矩阵（Assumption 3 满秩）。

\end{document}
